% Chapter on Terraform Infrastructure Management
% Add this as % Chapter on Terraform Infrastructure Management
% Add this as % Chapter on Terraform Infrastructure Management
% Add this as % Chapter on Terraform Infrastructure Management
% Add this as \input{sections/terraform-infra} in lab-manual.tex after the Ansible chapter

\chapter{Terraform Infrastructure Management}

Terraform configurations in this lab manage cloud infrastructure across two providers: DigitalOcean and Google Cloud Platform. All configurations are version-controlled and include automated security scanning via GitHub Actions.

\section{DigitalOcean Infrastructure}

The \texttt{terraform/digital\_ocean/} directory manages the external-facing web presence for bitsmasher.net.

\subsection{Provisioned Resources}

\begin{itemize}
    \item \textbf{Droplet:} 512MB Debian 12 in lon1 region (www.bitsmasher.net)
    \item \textbf{Domain:} bitsmasher.net with full DNS management
    \item \textbf{DNS Records:} A, MX, TXT (SPF/DMARC/DKIM), NS, and ACME challenge records
\end{itemize}

\subsection{Key DNS Records}

\begin{center}
\begin{tabular}{lll}
\hline
Name & Type & Purpose \\ \hline
www.bitsmasher.net & A & Web server host (178.62.60.55) \\
@ & MX & mail.protonmail.ch 10 (ProtonMail routing) \\
@ & TXT & SPF and DMARC policies \\
protonmail.\_domainkey & CNAME & DKIM email verification \\
\_acme-challenge.*.lab & TXT & Let's Encrypt certificate challenges \\ \hline
\end{tabular}
\end{center}

\subsection{Configuration Notes}

The droplet includes:
\begin{itemize}
    \item Automated backups enabled (immutable against accidental deletion)
    \item SSH key injection via terraform data source
    \item No external monitoring (consider enabling DO agent for uptime tracking)
\end{itemize}

Provider authentication uses a DigitalOcean API token managed through pass/direnv integration. The token is passed at runtime, never committed to the repository.

\section{Google Cloud Platform Infrastructure}

The \texttt{terraform/google/} directory manages GCP resources for Stash House free-tier operations.

\subsection{Resources Managed}

\begin{itemize}
    \item \textbf{VPC Network:} \texttt{stash-house-vpc} dedicated network
    \item \textbf{Compute Instance:} e2-micro Debian 12 (Always Free Tier eligible) in us-central1-a
    \item \textbf{Firewall Rule:} \texttt{allow-ssh} opening port 22 to the VPC
    \item \textbf{Billing Budget Alert:} Tripwire alert at \$0.01 threshold for free-tier monitoring
\end{itemize}

The GCP configuration is intentionally minimal, designed around Google's Always Free Tier limits to maintain zero-cost infrastructure.

\section{Security Scanning}

All Terraform configurations are automatically scanned by two GitHub Actions workflows:
\begin{itemize}
    \item \texttt{tfsec-sarif.yml} — IaC security scanning (generates SARIF reports)
    \item \texttt{tfsec-comment.yml} — Posts security findings as PR comments
\end{itemize}

Additionally, the DigitalOcean directory includes existing \texttt{README.md} with doctl and pass integration instructions for manual operations.

\section{Testing}

Terraform configurations are validated through:
\begin{itemize}
    \item Go-based integration tests (\texttt{terraform/test/main\_test.go})
    \item Plan dry-run validation via GitHub Actions (markdown.yml workflow)
    \item Manual state reconciliation checks (compare terraform state against live infrastructure)
\end{itemize}

State drift detection is recommended as a regular maintenance task — compare \texttt{terraform state list} output against actual cloud resource inventories quarterly.
 in lab-manual.tex after the Ansible chapter

\chapter{Terraform Infrastructure Management}

Terraform configurations in this lab manage cloud infrastructure across two providers: DigitalOcean and Google Cloud Platform. All configurations are version-controlled and include automated security scanning via GitHub Actions.

\section{DigitalOcean Infrastructure}

The \texttt{terraform/digital\_ocean/} directory manages the external-facing web presence for bitsmasher.net.

\subsection{Provisioned Resources}

\begin{itemize}
    \item \textbf{Droplet:} 512MB Debian 12 in lon1 region (www.bitsmasher.net)
    \item \textbf{Domain:} bitsmasher.net with full DNS management
    \item \textbf{DNS Records:} A, MX, TXT (SPF/DMARC/DKIM), NS, and ACME challenge records
\end{itemize}

\subsection{Key DNS Records}

\begin{center}
\begin{tabular}{lll}
\hline
Name & Type & Purpose \\ \hline
www.bitsmasher.net & A & Web server host (178.62.60.55) \\
@ & MX & mail.protonmail.ch 10 (ProtonMail routing) \\
@ & TXT & SPF and DMARC policies \\
protonmail.\_domainkey & CNAME & DKIM email verification \\
\_acme-challenge.*.lab & TXT & Let's Encrypt certificate challenges \\ \hline
\end{tabular}
\end{center}

\subsection{Configuration Notes}

The droplet includes:
\begin{itemize}
    \item Automated backups enabled (immutable against accidental deletion)
    \item SSH key injection via terraform data source
    \item No external monitoring (consider enabling DO agent for uptime tracking)
\end{itemize}

Provider authentication uses a DigitalOcean API token managed through pass/direnv integration. The token is passed at runtime, never committed to the repository.

\section{Google Cloud Platform Infrastructure}

The \texttt{terraform/google/} directory manages GCP resources for Stash House free-tier operations.

\subsection{Resources Managed}

\begin{itemize}
    \item \textbf{VPC Network:} \texttt{stash-house-vpc} dedicated network
    \item \textbf{Compute Instance:} e2-micro Debian 12 (Always Free Tier eligible) in us-central1-a
    \item \textbf{Firewall Rule:} \texttt{allow-ssh} opening port 22 to the VPC
    \item \textbf{Billing Budget Alert:} Tripwire alert at \$0.01 threshold for free-tier monitoring
\end{itemize}

The GCP configuration is intentionally minimal, designed around Google's Always Free Tier limits to maintain zero-cost infrastructure.

\section{Security Scanning}

All Terraform configurations are automatically scanned by two GitHub Actions workflows:
\begin{itemize}
    \item \texttt{tfsec-sarif.yml} — IaC security scanning (generates SARIF reports)
    \item \texttt{tfsec-comment.yml} — Posts security findings as PR comments
\end{itemize}

Additionally, the DigitalOcean directory includes existing \texttt{README.md} with doctl and pass integration instructions for manual operations.

\section{Testing}

Terraform configurations are validated through:
\begin{itemize}
    \item Go-based integration tests (\texttt{terraform/test/main\_test.go})
    \item Plan dry-run validation via GitHub Actions (markdown.yml workflow)
    \item Manual state reconciliation checks (compare terraform state against live infrastructure)
\end{itemize}

State drift detection is recommended as a regular maintenance task — compare \texttt{terraform state list} output against actual cloud resource inventories quarterly.
 in lab-manual.tex after the Ansible chapter

\chapter{Terraform Infrastructure Management}

Terraform configurations in this lab manage cloud infrastructure across two providers: DigitalOcean and Google Cloud Platform. All configurations are version-controlled and include automated security scanning via GitHub Actions.

\section{DigitalOcean Infrastructure}

The \texttt{terraform/digital\_ocean/} directory manages the external-facing web presence for bitsmasher.net.

\subsection{Provisioned Resources}

\begin{itemize}
    \item \textbf{Droplet:} 512MB Debian 12 in lon1 region (www.bitsmasher.net)
    \item \textbf{Domain:} bitsmasher.net with full DNS management
    \item \textbf{DNS Records:} A, MX, TXT (SPF/DMARC/DKIM), NS, and ACME challenge records
\end{itemize}

\subsection{Key DNS Records}

\begin{center}
\begin{tabular}{lll}
\hline
Name & Type & Purpose \\ \hline
www.bitsmasher.net & A & Web server host (178.62.60.55) \\
@ & MX & mail.protonmail.ch 10 (ProtonMail routing) \\
@ & TXT & SPF and DMARC policies \\
protonmail.\_domainkey & CNAME & DKIM email verification \\
\_acme-challenge.*.lab & TXT & Let's Encrypt certificate challenges \\ \hline
\end{tabular}
\end{center}

\subsection{Configuration Notes}

The droplet includes:
\begin{itemize}
    \item Automated backups enabled (immutable against accidental deletion)
    \item SSH key injection via terraform data source
    \item No external monitoring (consider enabling DO agent for uptime tracking)
\end{itemize}

Provider authentication uses a DigitalOcean API token managed through pass/direnv integration. The token is passed at runtime, never committed to the repository.

\section{Google Cloud Platform Infrastructure}

The \texttt{terraform/google/} directory manages GCP resources for Stash House free-tier operations.

\subsection{Resources Managed}

\begin{itemize}
    \item \textbf{VPC Network:} \texttt{stash-house-vpc} dedicated network
    \item \textbf{Compute Instance:} e2-micro Debian 12 (Always Free Tier eligible) in us-central1-a
    \item \textbf{Firewall Rule:} \texttt{allow-ssh} opening port 22 to the VPC
    \item \textbf{Billing Budget Alert:} Tripwire alert at \$0.01 threshold for free-tier monitoring
\end{itemize}

The GCP configuration is intentionally minimal, designed around Google's Always Free Tier limits to maintain zero-cost infrastructure.

\section{Security Scanning}

All Terraform configurations are automatically scanned by two GitHub Actions workflows:
\begin{itemize}
    \item \texttt{tfsec-sarif.yml} — IaC security scanning (generates SARIF reports)
    \item \texttt{tfsec-comment.yml} — Posts security findings as PR comments
\end{itemize}

Additionally, the DigitalOcean directory includes existing \texttt{README.md} with doctl and pass integration instructions for manual operations.

\section{Testing}

Terraform configurations are validated through:
\begin{itemize}
    \item Go-based integration tests (\texttt{terraform/test/main\_test.go})
    \item Plan dry-run validation via GitHub Actions (markdown.yml workflow)
    \item Manual state reconciliation checks (compare terraform state against live infrastructure)
\end{itemize}

State drift detection is recommended as a regular maintenance task — compare \texttt{terraform state list} output against actual cloud resource inventories quarterly.
 in lab-manual.tex after the Ansible chapter

\chapter{Terraform Infrastructure Management}

Terraform configurations in this lab manage cloud infrastructure across two providers: DigitalOcean and Google Cloud Platform. All configurations are version-controlled and include automated security scanning via GitHub Actions.

\section{DigitalOcean Infrastructure}

The \texttt{terraform/digital\_ocean/} directory manages the external-facing web presence for bitsmasher.net.

\subsection{Provisioned Resources}

\begin{itemize}
    \item \textbf{Droplet:} 512MB Debian 12 in lon1 region (www.bitsmasher.net)
    \item \textbf{Domain:} bitsmasher.net with full DNS management
    \item \textbf{DNS Records:} A, MX, TXT (SPF/DMARC/DKIM), NS, and ACME challenge records
\end{itemize}

\subsection{Key DNS Records}

\begin{center}
\begin{tabular}{lll}
\hline
Name & Type & Purpose \\ \hline
www.bitsmasher.net & A & Web server host (178.62.60.55) \\
@ & MX & mail.protonmail.ch 10 (ProtonMail routing) \\
@ & TXT & SPF and DMARC policies \\
protonmail.\_domainkey & CNAME & DKIM email verification \\
\_acme-challenge.*.lab & TXT & Let's Encrypt certificate challenges \\ \hline
\end{tabular}
\end{center}

\subsection{Configuration Notes}

The droplet includes:
\begin{itemize}
    \item Automated backups enabled (immutable against accidental deletion)
    \item SSH key injection via terraform data source
    \item No external monitoring (consider enabling DO agent for uptime tracking)
\end{itemize}

Provider authentication uses a DigitalOcean API token managed through pass/direnv integration. The token is passed at runtime, never committed to the repository.

\section{Google Cloud Platform Infrastructure}

The \texttt{terraform/google/} directory manages GCP resources for Stash House free-tier operations.

\subsection{Resources Managed}

\begin{itemize}
    \item \textbf{VPC Network:} \texttt{stash-house-vpc} dedicated network
    \item \textbf{Compute Instance:} e2-micro Debian 12 (Always Free Tier eligible) in us-central1-a
    \item \textbf{Firewall Rule:} \texttt{allow-ssh} opening port 22 to the VPC
    \item \textbf{Billing Budget Alert:} Tripwire alert at \$0.01 threshold for free-tier monitoring
\end{itemize}

The GCP configuration is intentionally minimal, designed around Google's Always Free Tier limits to maintain zero-cost infrastructure.

\section{Security Scanning}

All Terraform configurations are automatically scanned by two GitHub Actions workflows:
\begin{itemize}
    \item \texttt{tfsec-sarif.yml} — IaC security scanning (generates SARIF reports)
    \item \texttt{tfsec-comment.yml} — Posts security findings as PR comments
\end{itemize}

Additionally, the DigitalOcean directory includes existing \texttt{README.md} with doctl and pass integration instructions for manual operations.

\section{Testing}

Terraform configurations are validated through:
\begin{itemize}
    \item Go-based integration tests (\texttt{terraform/test/main\_test.go})
    \item Plan dry-run validation via GitHub Actions (markdown.yml workflow)
    \item Manual state reconciliation checks (compare terraform state against live infrastructure)
\end{itemize}

State drift detection is recommended as a regular maintenance task — compare \texttt{terraform state list} output against actual cloud resource inventories quarterly.
